\subsection{The literal mod-four pure sphere quotient and its additional fusion correction}

\begin{result}{A.4.31}{Proved}{The mod-eight crossing action admits an explicit pure correction at multiplier five}
\label{thm:sphere10-quaternary-fusion}
Retain the actual marked ten-point sphere, its mapping-class
groups $P\triangleleft M$, and its geometric Klein four group
$H$ from A.4.29.  Define the actual quotient
\[
 E_4=M/([P,P]P^4),\qquad V_4=P/([P,P]P^4).
                                                               \tag{QFU.1}
\]
Then $V_4\cong(\mathbb Z/4)^{35}$ and
$1\to V_4\to E_4\to\mathfrak S_{10}\to1$ is exact.
In the named crossing coordinates below, evaluation of the
standard braid formulas at any odd multiplier gives the genuine
automorphism
\[
 \Theta_\lambda(x,\pi)
   =\left(\lambda x+\frac{\lambda-1}{2}\,
                  \overline{c_\pi},\pi\right),
       \qquad \lambda\in(\mathbb Z/8)^\times.
                                                               \tag{QFU.2}
\]
The division is performed on an odd integral representative of
$\lambda$ before reduction modulo four.  The assignment depends
only on $\lambda\pmod8$ and obeys
$\Theta_\lambda\Theta_\mu=\Theta_{\lambda\mu}$.

The literal pure braid $W$ specified in (QFU.8)--(QFU.9)
has the following stronger correction property:
\[
     \operatorname{Inn}(W)\Theta_5
       \text{ fixes pointwise the entire inverse image of }
          \langle a,b\rangle\text{ in }E_4,
                                                               \tag{QFU.3}
\]
where $a,b$ are the two labelled permutations of $H$.
In particular the same correction works for every section
over these permutations.  For the specified geometric $H$,
all four automorphisms in (QFU.2) admit actual geometric
corrections fixing $H$ pointwise.  Their correction maps are
listed explicitly in (QFU.11).  Reduction $E_4\to E_2$
commutes with every standard action.
\end{result}

\begin{proof}
\par\medskip\noindent\textbf{The integral relation lattice and the actual quotient.}\quad
Use the $45$ edge generators $e_{ij}$ with
$0\leq i<j\leq9$, and write
\[
 L_{\mathbb Z}=\langle s_0,\ldots,s_9,t\rangle
       \subset\mathbb Z^{45},\qquad
 s_i=\sum_{j\ne i}e_{ij},\quad t=\sum_{i<j}e_{ij}.
                                                               \tag{QFU.4}
\]
The actual pure disk braid abelianization and the two
sphere mapping-class relations give
$P^{\rm ab}=\mathbb Z^{45}/L_{\mathbb Z}$, by the
source calculation in A.4.29.  The same exact presentation is
\cite[Theorem~3.2, printed p.~6]{CA_EndoKamada2016}.
Here the relation lattice is saturated over the integers,
not merely of the expected rank over a residue field.
Indeed $2t=\sum_{i=0}^9s_i$, so
$s_0,\ldots,s_8,t$ generate $L_{\mathbb Z}$.
On the ten columns
$e_{09},e_{19},\ldots,e_{89},e_{01}$, these ten rows
have the minor
\[
 \begin{pmatrix}I_9&v\\ \boldsymbol1^T&1\end{pmatrix},
 \qquad v=(1,1,0,\ldots,0)^T,
 \qquad\det=1-\boldsymbol1^Tv=-1.
                                                               \tag{QFU.5}
\]
They are consequently an integral basis of a direct summand.
It follows that $P^{\rm ab}\cong\mathbb Z^{35}$ and
that its reduction modulo four is exactly $V_4$ in
(QFU.1).  In particular every vector of this module is the
image of an actual pure braid.

For a permutation $\pi$, let $c_\pi$ be the integral
inversion vector defined in A.4.29: its $\{i,j\}$ coordinate
is one precisely when $\pi^{-1}(i)>\pi^{-1}(j)$ for
$i<j$.  Actual braid crossings modulo eight have the form
$c_\pi+2x$.  Quotienting the pure coordinate $x$ by
$L_{\mathbb Z}$ and by four gives the complete model
\[
\begin{gathered}
 (x,\pi)(y,\tau)
    =(x+\pi y+\kappa_4(\pi,\tau),\pi\tau),\\
 \kappa_4(\pi,\tau)
    =\overline{\frac{c_\pi+\pi c_\tau-c_{\pi\tau}}2}
       \quad\text{in }V_4.
\end{gathered}                                                  \tag{QFU.6}
\]
The numerator is even coordinatewise.  The cocycle equation
follows by cancellation before reducing coefficients.
This model is the actual quotient of $M$: on its pure
subgroup it is exactly the just-computed mod-four
abelianization, and its permutation map is unchanged.
Thus neither the extension nor its module is inferred from
an order comparison alone.

\par\medskip\noindent\textbf{Every odd standard action and its source relations.}\quad
The two arguments $y_i$ and $\sigma_i^2$ of each standard
braid conjugator both lie in $P$.  Their images in $V_4$
commute, so every commutator proword in those arguments
has value one.  The standard assignment is therefore
$\sigma_i\mapsto\sigma_i^\lambda$.
The square of $\sigma_i$ is the primitive edge class
$\overline e_{i-1,i}$ of order four, so $\sigma_i$
has order eight in $E_4$.

On an actual crossing vector, replacing each generator
by its odd $\lambda$th power preserves its adjacent
transposition and multiplies its signed crossing coordinate
by $\lambda$.  Concatenation uses the same permutation
action on edges, so this replacement multiplies the entire
crossing vector by $\lambda$.  The subgroup $2L_{\mathbb Z}$
of crossing relations is preserved.  Hence this assignment
descends to an endomorphism of the exact model (QFU.6),
giving (QFU.2).  Multiplication of crossing coordinates
proves the composition rule, and multiplication by the
inverse unit proves that it is an automorphism.
This also proves preservation of all source relations:
the braid relations have equal signed crossing vectors
and permutations, and the two sphere relations have
crossing vectors in $2L_{\mathbb Z}$.  The retained
program separately substitutes all four units into the
two literal sphere words as a check on this source map.

The action on $V_4$ is multiplication by $\lambda\pmod4$.
Writing $\beta_4(\pi)=\overline{c_\pi}$, the two
actions needed for fusion are
\[
 \Theta_5(x,\pi)=(x+2\beta_4(\pi),\pi),\qquad
 \Theta_7(x,\pi)=(-x-\beta_4(\pi),\pi)=\jmath(x,\pi).
                                                               \tag{QFU.7}
\]
The second is the actual braid mirror.  These formulas
also give the exact reduction to (BFU.2), and
$\Theta_3=\Theta_7\Theta_5$.

\par\medskip\noindent\textbf{A pure correction independent of the section.}\quad
Use the permutations
$a=(2\ 3)(4\ 5)(6\ 7)(8\ 9)$ and
$b=(0\ 1)(4\ 5)(6\ 8)(7\ 9)$.
In the edge quotient over $\mathbb Z/4$ define
\[
\begin{aligned}
 w={}&3e_{02}+2e_{03}+2e_{04}+e_{05}+e_{06}+3e_{08}\\
    &+3e_{12}+3e_{14}+3e_{16}+e_{18}
       +2e_{27}+2e_{28}\\
    &+2e_{34}+2e_{48}+2e_{56}+2e_{67}+2e_{68}.
\end{aligned}                                                   \tag{QFU.8}
\]
Let $w_{ij}\in\{0,1,2,3\}$ be the displayed coefficient,
and specify the actual pure word, in increasing lexicographic
edge order, by
\[
\begin{gathered}
 A_{ij}=\sigma_j\cdots\sigma_{i+2}\sigma_{i+1}^{\,2}
                  \sigma_{i+2}^{-1}\cdots\sigma_j^{-1},
         \qquad 0\leq i<j\leq9,\\
 W=\prod_{i<j}^{\rm lex}A_{ij}^{\,w_{ij}}.
\end{gathered}                                                  \tag{QFU.9}
\]
An empty product on either side of the square is omitted.
Thus $A_{i,i+1}=\sigma_{i+1}^2$.
Each $A_{ij}$ has pure linking coordinate $e_{ij}$, so
$W$ maps to $w$.  Its displayed expanded word has length
$258$; the program checks its actual crossings modulo eight.

Here are complete relation witnesses for the required
two identities.  Interpret $w$ first as the displayed
vector in the full edge module $(\mathbb Z/4)^{45}$.
Direct permutation of its coordinates gives
\[
\begin{aligned}
 w-aw-2c_a
   &=s_0+3s_1+2s_3+2s_5+2s_7+2s_8,\\
 w-bw-2c_b
   &=2s_0+2s_1+2s_3+2s_4+2s_5+2t.
\end{aligned}                                                   \tag{QFU.10}
\]
These equalities hold in the full $45$-coordinate module
modulo four and can be checked edge by edge; the two
inversion vectors are exactly those listed in (BFU.13).
After dividing by (QFU.4), they prove
$(1-a)w=2\beta_4(a)$ and
$(1-b)w=2\beta_4(b)$.

Conjugation by the pure element $w$ sends every
$(x,\pi)$ to $(x+(1-\pi)w,\pi)$.
Combining this formula with (QFU.7) and (QFU.10)
shows that $\operatorname{Inn}(w)\Theta_5$ is the identity
on each of the entire fibres over $a$ and $b$; it is
already the identity on $V_4$.
Those fibres and $V_4$ generate the inverse image of
$\langle a,b\rangle$.  The automorphism therefore fixes
that full inverse image pointwise, proving (QFU.3).
No lift of a binary linear solution is used: the displayed
relations are solved and verified over $\mathbb Z/4$ itself.

\par\medskip\noindent\textbf{The geometric mirror and all four corrections.}\quad
Choose the ordinary reflection $J$ defining the standard
braid marking.  It is an orientation-reversing involution
fixing the ten labels individually, and its conjugation on
$M$ is the standard mirror $\jmath$.
On the actual configuration (BFU.4), complex conjugation
$R:z\mapsto\overline z$ is another orientation-reversing
involution.  It interchanges precisely the two labels
of $C_0$, and it commutes with both actual maps
$z\mapsto-z$ and $z\mapsto1/z$.
Consequently $g=JR$ is an actual orientation-preserving
mapping class with permutation $c=(4\ 5)$ and
$\jmath(h)=ghg^{-1}$ for every $h\in H$.
This is an equality of the named source mapping classes,
so it remains true in $E_4$.

Together with the explicit word $W$, the corrections are
\[
\begin{array}{c|cccc}
 \lambda\pmod8&1&5&7&3\\ \hline
 k_\lambda&1&W&g^{-1}&g^{-1}\jmath(W)
\end{array}
\qquad
 \operatorname{Inn}(k_\lambda)\Theta_\lambda
                       \big|_H=1.
                                                               \tag{QFU.11}
\]
For the last column, compose
$\operatorname{Inn}(g^{-1})\jmath$ with
$\operatorname{Inn}(W)\Theta_5$; their product is exactly
$\operatorname{Inn}(g^{-1}\jmath(W))\Theta_3$.
The corrections for $1$ and $5$ are pure.  Those for $7$
and $3$ have permutation $c$, and hence are an actual
pure mapping class times the half twist $\sigma_5$.
This identifies the geometric correction and its labelled
source map without asserting that the binary correction
word itself survives unchanged at the next modulus.

The complete certificate is
\path{work/ghost_code/check_sphere10_quaternary_extension.py}.
It retains the saturated integral minor, the exact
mod-four quotient matrices, the colored extension cocycle,
the full correction system and its literal pure braid.
Its elimination has $23$ unit pivots and one residual
two-pivot, followed by direct checks of (QFU.10) in all
$45$ edge coordinates and of (QFU.3) on all four
permutation fibres.  It also checks the reduction to the
binary action and both global sphere relations for all
four units.  The text output, JSON and matrix archive are
retained beside the source.

This supplies the exact mod-four pure base and its standard
action for the relative ternary construction.  The map of
that larger construction to the capped mixed target and
the corresponding finite pentagon comparison are not
asserted here.
\end{proof}
